\documentclass[11pt,a4paper]{article}
\usepackage[T1]{fontenc}
\usepackage[utf8]{inputenc}
\usepackage{lmodern,microtype}
\usepackage[margin=.85in]{geometry}
\usepackage{amsmath,amssymb,amsthm,mathtools,aliascnt}
\usepackage{booktabs,tabularx,array,enumitem,needspace,float}
\usepackage[section]{placeins}
\usepackage{xcolor,tikz,mdframed}
\usetikzlibrary{arrows.meta,positioning,calc}
\usepackage[hidelinks]{hyperref}
\usepackage[nameinlink,noabbrev]{cleveref}
\usepackage{fancyhdr}
\hypersetup{pdftitle={Sync HotStuff RR v3.1},pdfsubject={Proposal timers and terminal-QC handoff}}
\setlength{\headheight}{14pt}
\pagestyle{fancy}\fancyhf{}
\lhead{\small Sync HotStuff RR v3.1}\rhead{\small L1-scheduled leaders}
\cfoot{\thepage}
\setlength{\parindent}{0pt}\setlength{\parskip}{4pt}
\setlength{\emergencystretch}{2em}
\setlist{nosep,leftmargin=*}
\newcolumntype{Y}{>{\raggedright\arraybackslash}X}
\definecolor{rrRed}{HTML}{A51C30}
\definecolor{rrBlue}{HTML}{2563EB}
\definecolor{rrTeal}{HTML}{0F766E}
\newcommand{\RR}[1]{{\color{rrRed}#1}}
\newcommand{\QC}{\mathcal C}
\newcommand{\Hash}{\mathsf H}
\newcommand{\rank}{\operatorname{rank}}
\newcommand{\tagg}[1]{\ensuremath{\mathsf{#1}}}
\newtheorem{definition}{Definition}[section]
\newaliascnt{assumption}{definition}
\newtheorem{assumption}[assumption]{Assumption}
\aliascntresetthe{assumption}
\newaliascnt{lemma}{definition}
\newtheorem{lemma}[lemma]{Lemma}
\aliascntresetthe{lemma}
\newaliascnt{theorem}{definition}
\newtheorem{theorem}[theorem]{Theorem}
\aliascntresetthe{theorem}
\newaliascnt{proposition}{definition}
\newtheorem{proposition}[proposition]{Proposition}
\aliascntresetthe{proposition}
\crefname{assumption}{assumption}{assumptions}
\Crefname{assumption}{Assumption}{Assumptions}
\crefname{lemma}{lemma}{lemmas}
\Crefname{lemma}{Lemma}{Lemmas}
\crefname{theorem}{theorem}{theorems}
\Crefname{theorem}{Theorem}{Theorems}
\newcounter{protocol}
\newenvironment{protocol}[2]{\refstepcounter{protocol}\label{#1}
\begin{mdframed}[nobreak=true,linewidth=.6pt,skipabove=8pt,skipbelow=8pt,
innertopmargin=7pt,innerbottommargin=7pt]\small
\textbf{Protocol \theprotocol. #2}\par\smallskip}{\end{mdframed}}
\crefname{protocol}{Protocol}{Protocols}
\Crefname{protocol}{Protocol}{Protocols}
\title{Sync HotStuff RR v3.1\\\large Proposal Timers and Terminal-QC Handoff}
\author{}\date{12 October 2026}
\begin{document}
\maketitle
\begin{abstract}
An L1 authority assigns the leader of each generation and term.
A term certifies $K\ge1$ fresh blocks; its terminal certificate admits
the next term without waiting for commitment. Replicas forward that
certificate to the next leader as a backup to ordinary vote delivery.
Every accepted DATA proposal starts one local $2\Delta$ timer, including
a previously unseen proposal from an earlier term of the same generation.
No certificate starts a separate timer. A complete quit proof freezes the
generation and invokes Sync HotStuff's lock-based recovery into $(g+1,0)$.
We prove safety under authenticated synchrony and eventual liveness under
explicit L1 availability, permanent equivocation exclusion, and fairness
assumptions. We also give counterexamples to earlier-term vote skipping,
shortened progress deadlines, and unconditional liveness without exclusion.
\end{abstract}

\section{Purpose and replication guarantees}
\label{sec:scope}
This specification revises the standard-synchronous Sync HotStuff
protocol~\cite{sync}. It separates ordinary leader rotation from recovery:
\[
 (g,t)\xrightarrow{\text{terminal DATA QC}}(g,t+1),
 \qquad
 g\xrightarrow{\text{complete quit proof}}(g+1,0).
\]
Ordinary rotation preserves commit timers. Recovery cancels all unfinished
timers in the departing generation. Compared with v3, this version removes
QC-start timers and vote skipping, gives explicit parent-selection rules,
and adds terminal-QC forwarding to the next leader. Both protocols use the
same local term-entry event to start their progress checks.

\begin{definition}[State machine replication]
Blocks order commands for a deterministic state machine. Committing a
block commits its uncommitted ancestors in order. \emph{Safety} requires
correct replicas' committed logs to be prefix-comparable and never revoked.
\emph{Validity} requires application-valid committed blocks.
\emph{Request liveness} requires every correct-client request that remains
admissible until ordered to be committed eventually by every correct replica.
\end{definition}
We first prove unbounded committed-height progress, then derive request
liveness using an explicit input-fairness assumption. Normally completed
terms contain $K$ fresh blocks. A faulty term can produce none and be
replaced. Terms are defined by certified block counts, not fixed physical
durations. Replicas use a global clock but may enter a term at different times.

\section{Model and service assumptions}
\label{sec:model}
There are $n=2f+1$ fixed voting members, at most $f$ permanently Byzantine.
Initialization uses a common genesis and initial authority; correct
replicas enter $(0,0)$ within $\Delta$ of one another.

\Needspace{9\baselineskip}
\begin{assumption}[Authenticated synchronous communication and global clock]
\label{ass:network}
\leavevmode\par
\begin{itemize}
    \item Messages use authenticated communication channels.
    \item The actual delay of a correct-to-correct message is $\delta\le\Delta$;
    $\delta$ may differ across messages.
    \item All replicas share a global clock.
\end{itemize}
\end{assumption}
The delay bound includes delivery and processing of complete validation
dependencies. Each replica records its own entry and voting events on
the global clock.

\begin{assumption}[Finalized authority]
\label{ass:authority}
L1 is a public bulletin board. For each generation $g$, it supplies one
immutable authority $\tagg{AUTH}_g$ specifying a unique leader $L(g,t)$
from the fixed committee for every term $t$. At the finite global time
$d(g)$, that authority is finalized and simultaneously available to all
correct replicas, including frozen ones. L1 progresses independently
of L2 commitment.
\end{assumption}
An authority may be finalized in advance. The L1 epoch and recovery
generation need not coincide.

\begin{assumption}[Accountability]
\label{ass:exclusion}
Every admissible leader-equivocation proof revealed to a correct replica
is eventually submitted to L1 and processed. For each exposed identity $p$,
let $g_{\mathrm{bad}}$ be the generation identified by the first such proof
revealed to a correct replica. This is the generation of the proved
equivocation, not necessarily the generation current at revelation.
There is a finite removal generation $G_p>g_{\mathrm{bad}}$ such that
\[
 L(g,t)\ne p\quad\text{for every }g\ge G_p\text{ and every }t\ge0.
\]
Exclusion is permanent. The finite Byzantine identity set cannot be
replenished by new corruptions, joins, or replacement identities.
\end{assumption}
Both $g_{\mathrm{bad}}$ and $G_p$ are generation indices. No fixed bound
on $G_p-g_{\mathrm{bad}}$ is assumed. Exclusion takes effect in future
authorities without rewriting a finalized assignment.
This is an enforced eligibility property, not an economic claim that a
deposit deters faults. Timeout BLAME alone is not equivocation evidence.
Excluded proposers may remain faulty voters within the same $f$ bound.
Changing voting membership or quorum weights needs a separate
cross-configuration argument.

\begin{assumption}[Schedule and input fairness]
\label{ass:fairness}
For every $g_0$, some finite $g\ge g_0$ has correct $L(g,0)$.
For every fixed $g$ and $t_0$, some finite $t\ge t_0$ has correct $L(g,t)$.
Correct leaders can propose valid empty blocks.

Correct-client requests eventually reach all correct replicas. They have
a common admission order with finitely many predecessors per request, and
each request fits in a batch. Correct leaders include the oldest pending
admissible requests that fit. A request leaves the pending set only after
commitment; inclusion in the proposed ancestry avoids duplication but
does not by itself remove the request. Correct replicas retain and serve
validation data.
\end{assumption}
Safety uses \cref{ass:network,ass:authority}. Exclusion and fairness are
additional assumptions for eventual progress, not premises about a
Byzantine leader's willingness to produce blocks.

\section{Blocks, certificates, and message handling}
\label{sec:definitions}
\paragraph{Views and positions.}
A view is $(g,t)$, ordered lexicographically. Generation $g$ has an
opening block $A_g$; $A_0$ is genesis. Fresh position $j\in\{0,\ldots,K-1\}$
of term $(g,t)$ has height
\[
 k=h(A_g)+tK+j+1.
\]
Its block is $B_k=(x_k,\Hash(B_{\mathrm{parent}}))$, with
$h(B_{\mathrm{parent}})=k-1$. The term's terminal block has $j=K-1$.

\paragraph{Signed context and certificates.}
All signatures bind the protocol and voting configuration, L1 authority,
generation opening, message type, view, and the block's height and hash;
DATA signatures also bind position $j$. We omit these fixed context fields
from message displays. A DATA QC $\QC_{g,t}(B_k)$ contains $f+1$ matching
votes from distinct members. ANCHOR votes are separately tagged and
re-certify an unchanged opening block in $(g,0)$. A correct replica casts
at most one DATA vote per $(g,t,j)$ and one ANCHOR vote per generation.

Certificate rank is $(g,t,k)$. Equal-rank choices use a fixed total
tie-break on certified content, independently of signer subsets or
encodings. Genesis has a synthetic certificate below all ordinary ranks.
An opening's identity is $(g,h(A_g),\Hash(A_g))$.

\paragraph{Admission and complete validation.}
For $g>0$, the opening is a signed NEWVIEW from $L(g,0)$, justified by
a valid QC from a strictly older generation and an entry proof for $g$.
Its ANCHOR QC admits fresh DATA proposals. Genesis needs no ANCHOR exchange.
Term $t>0$ additionally requires all preceding terminal QCs in that
generation. Term zero is admitted by its generation-entry proof and
authority even before its leader sends NEWVIEW.

Validation checks signatures, authority, the opening and term admission,
position and height, the parent rule in \cref{tab:parent}, parent QCs,
hash links, and application validity throughout the required ancestry.
Leader authorization is checked against the public L1 authority.
Every relevant ancestor includes its signed proposal and block data;
a chain of certificate hashes alone is insufficient. Dependencies can
come from the message or a validated local cache. Correct senders provide
complete validation bundles; references are allowed only when the receiver
can resolve them within the same complete-message bound. Bundles list
each dependency once rather than recursively copying nested bundles.

\paragraph{Evidence takes priority.}
Before validating an enclosing message or discarding a duplicate, inspect
available signed material for an independently verifiable quit proof.
Evidence must itself have valid signatures and admission witnesses.
Missing or invalid outer data neither suppresses a complete proof nor
turns an unauthenticated conflict into one. Compare incoming histories
with retained histories, not just with other objects inside the message.
At an equal timestamp, process received messages before timer callbacks.

\paragraph{One proposal-processing rule.}
Every previously unvoted, admitted current-generation DATA proposal invokes
the Vote step, including proposals in attached histories, proposals from
earlier terms, and proposals whose blocks have already been committed as
ancestors. Process available ancestors before their children and before
terminal-QC entry. An earlier term does not make a proposal ineligible.
Deduplicate proposals by $(g,t,j,\Hash(B_k))$ after evidence inspection.
Different certificate representations do not create a new proposal.
Old-generation DATA proposals are retained but do not receive new votes
or timers. Future-generation messages wait for authenticated entry.

\paragraph{Active and frozen generations.}
New proposals, DATA/ANCHOR votes, and ordinary term transitions require
the current generation to be \tagg{Unfrozen}. A frozen generation still
accepts evidence, old votes, and QCs, observes the L1 bulletin board, and
executes recovery. No complete commitment is withdrawn.

\begin{table}[htbp]\centering\small
\caption{Parent selection for the next fresh position $j$ of term $(g,t)$.}
\label{tab:parent}
\begin{tabularx}{\linewidth}{@{}l Y Y@{}}
\toprule
Condition & $B_{\mathrm{parent}}$ & $Q_{\mathrm{parent}}$\\
\midrule
$j=0,t=0$ & Opening block $A_g$ & Its ANCHOR QC in $(g,0)$\\
$j=0,t>0$ & Position $K-1$ of $(g,t-1)$ & Its DATA QC in $(g,t-1)$\\
$1\le j<K$ & Position $j-1$ of $(g,t)$ & Its DATA QC in $(g,t)$\\
\bottomrule
\end{tabularx}
\smallskip
For $(g,t,j)=(0,0,0)$, use genesis and its initial certificate instead.
\end{table}

\section{Protocols}
\label{sec:protocols}
The shared definitions above apply to both boxes. Red text marks changes
or explicit extensions relative to standard-synchronous SH; explanations
and validation requirements are kept outside the execution steps.

\begin{protocol}{alg:steady}{Propose, vote, commit, and normal rotation}
\begin{enumerate}[label=\textbf{\arabic*)},itemsep=6pt]
\item \textbf{Propose.} If $r=L(g,t)$ in its current term, obtain the
parent and certificate for its next fresh position $j$
\RR{using \cref{tab:parent}}. Immediately broadcast
\[
 P=\langle\tagg{PROPOSE},g,t,j,B_k,Q_{\mathrm{parent}}\rangle_{L(g,t)}
\]
with the validation dependencies, where
$B_k=(x_k,\Hash(B_{\mathrm{parent}}))$.
Process this proposal locally under Vote. The parent need not be committed.

\item \textbf{Vote.} On a proposal $P$, \RR{process quit evidence first;
return if it freezes the current generation}. Then check completeness and validity. If data are
missing, retain $P$ and request them without voting; reject invalid $P$.
If already voted for $P$, stop. Otherwise forward $P$ with its dependencies,
broadcast
$\langle\tagg{DATA\mbox{-}VOTE},g,t,j,\Hash(B_k)\rangle_r$,
and start its commit timer at the local acceptance time $\tau_r(P)$,
expiring at $\tau_r(P)+2\Delta$.
\RR{The same rule processes unseen earlier-term proposals and attached
DATA history. A duplicate never restarts the timer.}

\item \textbf{Commit.} On expiry of $P$'s timer, commit $B_k$ and its
uncommitted ancestors. \RR{Normal rotation preserves timers; Quit cancels
unfinished generation timers. A QC alone starts no timer.}

\item \textbf{Terminal QC and rotation.} \RR{On first obtaining a valid
terminal DATA QC of $(g,t)$, process its unseen DATA history under Vote.
If still unfrozen, send the QC and its complete validation bundle to
$L(g,t+1)$ once. Enter $(g,t+1)$ if not already there or beyond it;
record this first local entry as $e_r(g,t+1)$ and disable completed-term
progress checks. Process a chain of terminal QCs in order. Old or duplicate
QCs never move the term backward or reset an entry time. The successor's
first proposal carries its parent's terminal QC.}
\end{enumerate}
\end{protocol}

The acceptance time $\tau_r(P)$ is when complete validation and the first
vote occur, not an earlier arrival time of an incomplete packet. A local
term starts at its first valid admission event, not at receipt of the
new leader's first proposal. Thus silence by a newly admitted leader does
not prevent its progress clock from starting.

\paragraph{Complete quit proofs.}
For an admitted generation $g$, a proof consists of either:
\begin{enumerate}[label=(\alph*)]
\item $f+1$ distinct signatures on the same $\langle\tagg{BLAME},g,t\rangle$
for an admitted term;
\item two different admissible DATA blocks signed by $L(g,t)$ at the same
$(g,t,j)$, with their admission histories; or
\item two valid NEWVIEWs signed by $L(g,0)$ selecting different opening
blocks for generation $g$.
\end{enumerate}
Different openings count even when their blocks are comparable, because
they assign different quota origins. Different QCs for the same opening
block are not equivocation. Evidence from any admitted term can freeze
its generation, including a term not yet locally reached. Old-generation
proofs are archival and may be submitted to L1; future-generation proofs
are buffered until entry. Neither resets an unrelated recovery.

\begin{protocol}{alg:recovery}{Blame, quit, and recovery}
\begin{enumerate}[label=\textbf{\arabic*)},itemsep=6pt]
\item \textbf{Blame.} For each $p=1,\ldots,K$, check progress at
\[
 e_r(g,t)+(2p+4)\Delta.
\]
If $(g,t)$ remains current and unfrozen and fewer than $p$ distinct fresh
positions have triggered this replica's DATA votes, broadcast
$\langle\tagg{BLAME},g,t\rangle_r$ once for that term.
\RR{ANCHOR votes do not count.}

\item \textbf{Quit.} On the first complete proof for the current generation
$g$, \RR{freeze all of $g$: stop its progress checks, new proposals, votes,
and normal transitions; cancel all unfinished proposal timers}.
Broadcast the proof and record the local quit time $a_r$ once.
Retain commitments and continue processing old votes and QCs.
\RR{Submit revealed equivocation to L1; local recovery does not await
that submission's finalization.}

\item \textbf{Status and entry.} \RR{The target is $(g+1,0)$. At}
\[
 \RR{e_r(g+1,0)=\max\{a_r+\Delta,d(g+1)\},}
\]
\RR{read $\tagg{AUTH}_{g+1}$ from L1}. Fix the highest-ranked certificate $Q_r$
among the retained lock and known valid QCs from generations through $g$
as the entry lock. Enter \RR{$(g+1,0)$, mark it unfrozen, and start its
progress clock}. Send
$\langle\tagg{STATUS},g+1,0,Q_r,E\rangle_r$
with validation dependencies to \RR{$L(g+1,0)$}; $E$ is the quit proof.
\RR{Immediately process buffered proofs for the new generation.
Later QCs update retained evidence, not this fixed entry lock.}

\item \textbf{New-view.} The new leader waits $2\Delta$ from its own entry.
If the new generation remains current and unfrozen, choose the highest
valid known QC $Q^*$ from generations through $g$ and broadcast
\[
 \langle\tagg{NEWVIEW},g+1,0,A,Q^*,E\rangle_{L(g+1,0)},
\]
where $Q^*$ certifies $A$, with complete dependencies.
\RR{Choose only one opening block for this generation.}

\item \textbf{First vote.} On NEWVIEW, \RR{process verifiable quit evidence
first}, then validate its leader, entry proof, anchor, QC, and dependencies.
If the new generation is current and unfrozen, no ANCHOR vote was cast,
and $\rank(Q^*)\ge\rank(Q_r)$, forward NEWVIEW and broadcast
$\langle\tagg{ANCHOR\mbox{-}VOTE},g+1,0,\Hash(A)\rangle_r$.
\RR{Start no timer and consume no fresh position. On its ANCHOR QC, the
leader resumes \cref{alg:steady} at $j=0$.}
\end{enumerate}
\end{protocol}

The entry-lock test governs ANCHOR voting only. A subsequent DATA proposal
backed by a valid ANCHOR QC is processed under \cref{alg:steady}, even at
a replica that declined the original ANCHOR vote. Missing NEWVIEW
dependencies are requested and awaited without voting; invalid messages
are rejected. The complete-message bound applies to correct transmissions.

\section{Execution times and progress deadlines}
\label{sec:timing}
\begin{lemma}[Generation entry and freeze propagation]
\label{lem:entry}
Correct entry times to $(g,0)$ differ by at most $\Delta$. If the first
correct freeze of $g$ is at $a$, every correct replica freezes it by
$a+\Delta$. All recover to $(g+1,0)$, with entry skew at most $\Delta$.
\end{lemma}
\begin{proof}
The initial skew is assumed. A proof is broadcast on first freeze.
Any correct replica already in $g$ receives it by $a+\Delta$.
A replica still finishing prior recovery enters $g$ within $\Delta$
of its first correct entry, hence by $a+\Delta$, and processes the
buffered proof. No correct replica can complete recovery from $g$ before
the first correct freeze of $g$.

All replicas use the same public availability time $d(g+1)$. Hence
\[
\begin{aligned}
 |e_r(g+1,0)-e_s(g+1,0)|
 &=\big|\max\{a_r+\Delta,d(g+1)\}
          -\max\{a_s+\Delta,d(g+1)\}\big|\\
 &\le |a_r-a_s|\le\Delta.
\end{aligned}
\]
This gives the next generation's entry bound directly.
A buffered proof may freeze the next generation immediately after entry;
it does not erase that entry event or change its target.
\end{proof}

\begin{lemma}[Terminal-QC handoff]
\label{lem:handoff}
Let $q$ be the first time a correct replica in generation $g$ has verified
a complete terminal QC for $(g,t)$. If the successor is correct and the
generation stays unfrozen, it proposes its first child by $q+\Delta$,
and all correct replicas process that proposal by $q+2\Delta$.
Without assuming a correct successor, all correct replicas still obtain
the terminal QC by $q+2\Delta$ if the generation remains unfrozen.
\end{lemma}
\begin{proof}
The first correct QC holder sends the complete bundle to the successor.
It arrives by $q+\Delta$. If that leader is still finishing generation
entry, \cref{lem:entry} places its entry no later than $q+\Delta$.
It has the signed terminal block and its QC and immediately proposes
the child, unless it already did so. Its proposal reaches all by
$q+2\Delta$.

For the final claim, a correct terminal-QC signer voted at some $s\le q$
and broadcast its complete proposal. All correct replicas have entered
the generation and process it by $s+\Delta$. Earlier-term proposals
remain eligible, so all cast or have already cast that DATA vote.
Their $f+1$ or more correct votes give the terminal QC everywhere by
$s+2\Delta\le q+2\Delta$.
\end{proof}
Here $q$ is a correct replica's verified QC receipt, not an adversary's
private assembly time. Ordinary term-entry skew is at most $2\Delta$
in this unfrozen execution, whereas recovery-generation entry skew is
$\Delta$. A common clock does not require simultaneous term entry.

\begin{figure}[htbp]\centering
\begin{tikzpicture}[x=1cm,y=-1cm,font=\small,
 msg/.style={-{Stealth[length=2mm]},semithick,draw=rrBlue},
 lab/.style={fill=white,inner sep=3pt,align=center}]
\path[use as bounding box] (0,-.5) rectangle (14.8,5.0);
\node[align=center] at (1.9,0) {First correct holder\\of terminal QC};
\node[align=center] at (7.35,0) {Correct next leader\\$L(g,t+1)$};
\node[align=center] at (12.6,0) {Other correct\\replicas};
\foreach \x in {1.9,7.35,12.6} \draw[black!30,dashed] (\x,.6)--(\x,4.4);
\node[anchor=east] at (1.7,1.15) {$q$};
\draw[msg] (1.9,1.15)--node[lab,above,sloped]{terminal QC + complete history}(7.35,2.05);
\node[lab,anchor=west] at (7.5,2.05) {by $q+\Delta$: enter; propose};
\draw[msg] (7.35,2.65)--node[lab,above,sloped]{first DATA proposal + parent QC}(12.6,3.65);
\draw[msg] (7.35,2.65)--(1.9,3.65);
\node[lab] at (7.35,4.35) {By $q+2\Delta$: all have processed the new proposal.};
\end{tikzpicture}
\caption{Normal handoff while the generation remains unfrozen. The arrows
show causal stages, not mandatory waiting. Each replica records its own
first terminal-QC entry time. Backup forwarding introduces no vote phase.}
\label{fig:handoff}
\end{figure}

\begin{lemma}[Correct-leader progress]
\label{lem:progress}
While a correct leader's term remains active and no correct replica has
frozen the generation, its $p$th fresh position triggers every correct
replica's DATA vote by that replica's deadline
$e_r(g,t)+(2p+4)\Delta$, for $p=1,\ldots,K$.
Thus a correct replica sends no progress BLAME against such a leader
before the first correct freeze.
\end{lemma}
\begin{proof}
For recovered term zero, put $e=\min_r e_r(g,0)$. The correct leader enters
by $e+\Delta$ and finishes its wait by $e+3\Delta$. All correct STATUS
messages arrive within that wait, so its selected QC meets every entry
lock. NEWVIEW reaches all by $e+4\Delta$; ANCHOR votes reach the leader
by $e+5\Delta$; its first fresh proposal reaches all by $e+6\Delta$.
Each further proposal costs at most $2\Delta$: votes back to the leader,
then the next proposal to all. The $p$th vote therefore occurs by
\[
 e+6\Delta+2(p-1)\Delta=e+(2p+4)\Delta
 \le e_r(g,0)+(2p+4)\Delta.
\]
Genesis has no ANCHOR exchange and is faster.

For $t>0$, let $q$ be as in \cref{lem:handoff} for the preceding terminal
QC. The first proposal arrives by $q+2\Delta$ and the $p$th by
$q+2p\Delta$. Every local entry is at or after $q$, so these arrivals
precede the stated deadlines. History proposals are processed even after
ordinary advancement, and completed-term checks are disabled. At equality,
received proposals are processed before a progress callback.
\end{proof}

\begin{figure}[htbp]\centering
\begin{tikzpicture}[node distance=4mm,font=\small,
 event/.style={draw=rrBlue!65,rounded corners=2pt,align=left,
 text width=12.2cm,inner sep=5pt},
 arr/.style={-{Stealth[length=1.7mm]},draw=rrBlue}]
\node[event] (a) {$e$: earliest correct entry to recovered term zero};
\node[event,below=of a] (b) {$e+\Delta$: latest possible correct leader entry};
\node[event,below=of b] (c) {$e+3\Delta$: leader completes its own $2\Delta$ wait; sends NEWVIEW};
\node[event,below=of c] (d) {$e+4\Delta$: all correct replicas send ANCHOR votes};
\node[event,below=of d] (f) {$e+5\Delta$: leader has ANCHOR QC; sends first fresh DATA proposal};
\node[event,below=of f,draw=rrTeal] (h) {$e+6\Delta$: all correct replicas vote; first progress check is met};
\foreach \u/\v in {a/b,b/c,c/d,d/f,f/h} \draw[arr] (\u.south)--(\v.north);
\end{tikzpicture}
\caption{Worst-case recovery timeline for a correct initializer. Subsequent
fresh positions add $2\Delta$ each, giving $(2p+4)\Delta$. ANCHOR votes
are not fresh progress. Uninterrupted DATA timers commit the first block
at all correct replicas by $e+8\Delta$.}
\label{fig:recovery}
\end{figure}
\FloatBarrier

\section{Safety analysis}
\label{sec:safety}
\begin{lemma}[One generation, one chain or a proof]
\label{lem:fork}
Two admitted histories in the same generation have the same opening and
are prefix-comparable, or their signed material gives a valid quit proof.
\end{lemma}
\begin{proof}
Different openings give two conflicting signed NEWVIEWs. With one opening,
the height formula assigns each fresh height to exactly one $(t,j)$ and
leader. The first divergent fresh height then gives two different signed
DATA proposals at that position. The histories provide their admission
witnesses. An older-generation fork in an inherited justification is not
by itself a new-generation equivocation.
\end{proof}

\begin{lemma}[A completed proposal timer protects its block]
\label{lem:timer}
If a correct replica votes for $B$ at $s$ and commits it at $s+2\Delta$,
every correct replica knows its DATA QC by that expiry. No conflicting
QC in the same generation, nor an ANCHOR QC for a different opening,
can exist, including a QC assembled later from old votes.
\end{lemma}
\begin{proof}
A correct freeze by $s+\Delta$ would reach the committer by $s+2\Delta$
and cancel its timer. Thus none occurs that early. The complete proposal
sent at $s$ reaches every correct replica by $s+\Delta$; even a replica
still finishing recovery has entered this generation by then, by
\cref{lem:entry}. Any incompatible retained history would expose a proof
and freeze that receiver by $s+\Delta$, a contradiction. Therefore every
correct replica has already cast or now casts the same DATA vote. Normal
term advancement and prior indirect commitment do not suppress that vote.
Their at least $f+1$ multicast votes provide the QC everywhere by expiry.

Any conflicting QC has a correct signer. If it processed $B$ first, its
retained evidence prevents the conflicting vote. Otherwise its conflicting
vote and forwarded signed history preceded that processing, no later
than $s+\Delta$. That history reaches the committer by expiry and exposes
the fork, contradicting commitment. An ANCHOR vote for another opening
similarly forwards the conflicting NEWVIEW. This reasoning bounds the
correct vote's time, not the later time of QC assembly. Evidence-first
handling covers simultaneous arrival and expiry.
\end{proof}

\begin{lemma}[Recovery locks preserve a completed timer]
\label{lem:lock}
If a correct timer commits $B$ in generation $g$, every correct entry
lock in a later generation ranks at least as high as its DATA QC.
\end{lemma}
\begin{proof}
Let $s$ be its vote time and $a$ the first correct freeze of $g$.
The preceding cancellation argument gives $a>s+\Delta$. Every successor
entry satisfies
$e_r(g+1,0)\ge a_r+\Delta\ge a+\Delta>s+2\Delta$.
The protecting QC is therefore known before every lock snapshot.
Highest-QC selection and retained locks preserve the rank bound in all
later generations. Additional L1 waiting only delays these snapshots.
If the generation never freezes, there is no later entry to check.
\end{proof}

\begin{lemma}[Unique extensibility]
\label{lem:extension}
For a block $B$ committed by a correct timer, every admitted QC ranked
at least as high as its protecting QC certifies an extension of $B$.
\end{lemma}
\begin{proof}
In its generation, \cref{lem:timer} excludes alternative certified
openings and conflicting QCs. With a fixed opening, higher terms own
strictly higher fresh heights, so a sufficiently ranked QC cannot certify
a strict ancestor of $B$.

If a later-generation counterexample exists, choose one of least rank.
An ANCHOR QC has a correct voter whose older justification meets its
entry lock, hence the protecting rank by \cref{lem:lock}. That justification
certifies the same nonextending block at a strictly smaller rank, a
contradiction. A DATA QC has a lower-ranked parent QC in the same later
generation, including at a term boundary. Its parent cannot extend $B$
if the child does not, so this too is a smaller counterexample. The
lexicographic nonnegative ranks are well-ordered; content tie-breaks do
not affect these strict descents.
\end{proof}

\begin{theorem}[Safety and validity]
\label{thm:safety}
Under \cref{ass:network,ass:authority}, correct committed logs are
prefix-comparable, never revoked, and application-valid.
\end{theorem}
\begin{proof}
Order the protecting QCs of any two direct timer commitments by rank.
\Cref{lem:extension} makes the higher-ranked block extend the lower-ranked
one. Indirect ancestor commitments inherit comparability. Commitments
are never withdrawn by either transition. Complete validation establishes
application validity. Neither proposer exclusion nor schedule fairness
is used in this argument.
\end{proof}

\section{Liveness analysis}
\label{sec:liveness}
\begin{lemma}[One timer rule suffices]
\label{lem:window}
Suppose a correct replica votes for and forwards a DATA proposal at $s$.
If no correct replica freezes that generation through $s+3\Delta$, every
correct replica commits its block by $s+3\Delta$; the originating timer
commits at $s+2\Delta$.
\end{lemma}
\begin{proof}
Every correct replica has entered and processes the proposal by $s+\Delta$.
An incompatible history would force a freeze, contrary to the premise.
Thus every correct replica votes by that time or has already voted.
Every such vote starts the one proposal timer. Earlier terms and
QC-carrying history use exactly the same rule. All timers expire by
$s+3\Delta$ without cancellation; prior commitments already satisfy
the conclusion. No QC-start timer is needed.
\end{proof}

\begin{lemma}[Only finitely many generations expose equivocation]
\label{lem:finite}
Under \cref{ass:exclusion}, only finitely many generation indices are
named by admissible equivocation proofs revealed to correct replicas.
\end{lemma}
\begin{proof}
Each identified leader is Byzantine, since correct leaders sign one
opening and one DATA block per position. For each identity $p$ ever exposed,
its first revealed proof names $g_{\mathrm{bad}}$ and accountability gives
a finite $G_p>g_{\mathrm{bad}}$. Any admissible equivocation by $p$ must
belong to a generation $g<G_p$: at $g\ge G_p$, the public authority
never assigns $p$ a leader position. The union of these finite ranges
over the finite exposed identity set is finite. A proof revealed later
still names its original generation and cannot restart a completed recovery.
\end{proof}
This does not bound interruptions by $f$. An identity can cause several
before exclusion becomes effective.

\begin{lemma}[A correct initializer obtains a commit window]
\label{lem:initializer}
Consider a generation with correct $L(g,0)$ in which no admissible
equivocation proof is revealed. Its first fresh proposal and local DATA
vote occur at a finite time $s$, and every correct replica commits that block by
$s+3\Delta$, even if QC rotation reaches a Byzantine successor.
\end{lemma}
\begin{proof}
Let $e$ be earliest entry. Before the first fresh proposal, no later
term is admitted: its terminal-QC ancestry would require that proposal.
The earliest possible correct term-zero BLAME is at $e+6\Delta$.
With equivocation excluded, no quit proof can freeze this initial interval,
and the recovery envelope produces the first proposal by $e+5\Delta$.
Thus $s$ exists. Genesis initialization is faster.

Suppose the first correct freeze were at $a\le s+3\Delta$.
It must use BLAME. Before that first freeze, \cref{lem:progress} excludes
any correct term-zero BLAME. For a later term, a correct BLAME signer
entered at some $x\ge s$: its admission history includes the initializer's
first proposal, signed at $s$. The first complaint is no earlier than
$x+6\Delta\ge s+6\Delta$, even if Byzantine shares were prepared in
advance. Thus no complete BLAME proof can cause the supposed freeze.
\Cref{lem:window} completes the claim.
\end{proof}

\begin{theorem}[Eventual committed-height progress]
\label{thm:height}
Under \cref{ass:network,ass:authority,ass:exclusion,ass:fairness}, every
correct replica eventually commits arbitrarily large heights.
\end{theorem}
\begin{proof}
If infinitely many generations are entered, \cref{lem:finite} leaves a
suffix without revealed equivocation. Schedule fairness supplies infinitely
many correct initializers there. By \cref{lem:initializer}, each commits
its first fresh block everywhere. Its certified opening extends prior
committed blocks by locks and unique extensibility, so this fresh child
has strictly greater height. Authority availability prevents a permanent
wait for any one next generation.

If only finitely many generations are entered, the last one cannot freeze:
the finite drain interval and eventual next authority would permit another
entry. Suppose this unfrozen generation had a last term. Every observed
terminal QC spreads through its honest signer's proposal and the resulting
correct votes, by \cref{lem:handoff}. Thus all correct replicas eventually
reach that last common term. If any correct replica votes for its terminal
proposal, forwarding and correct votes produce the terminal QC everywhere
and admit another term. Otherwise every correct replica has fewer than
$K$ votes at its last progress check and broadcasts BLAME. Their matching
shares force a freeze. Both contradict a last term.

Hence certified terms and heights increase without bound. Each terminal
QC has a correct signer who forwards its complete proposal; all correct
replicas vote under the same rule and their timers commit, since this
generation never freezes. This proves progress at every correct replica.
\end{proof}

\begin{theorem}[Request liveness]
\label{thm:requests}
Under the same assumptions, every persistently admissible correct-client
request is eventually committed by all correct replicas.
\end{theorem}
\begin{proof}
In the infinite-generation case, infinitely many correct initializers
after the last equivocation generation commit their first fresh blocks
everywhere. In the finite-generation case, term fairness gives infinitely
many correct leaders in the last unfrozen generation; their fresh blocks
commit. Fix a request after dissemination. Only finitely many requests
precede it. Each such successful correct proposal includes the oldest
fitting pending requests or commits them as ancestors. Speculative
inclusion does not remove a pending request. Induction over the finite
predecessors therefore commits the chosen request, and safety preserves it.
\end{proof}

\section{Attacks on tempting simplifications and weaker assumptions}
\label{sec:attacks}
The preceding theorems concern the specified protocol and assumptions.
The following executions identify exactly what fails when a rule or
assumption is removed. They are included to distinguish a proved conditional
guarantee from an unconditional claim.

\paragraph{A. Skip old proposals but remove their QC timers.}
Take $n=3$, $f=1$, $K=1$, correct replicas $H_0,H_1$, and Byzantine $F$.
In each generation let $H_0$ initialize and $F$ follow it; later terms
can cycle fairly over all identities. $H_0$ proposes $B$ and starts its
timer. $F$ supplies its vote, forms a terminal QC, and gives that QC to
$H_1$ before $H_1$ processes the proposal. If $H_1$ first advances and
then refuses earlier-term proposals, it never starts $B$'s timer.
$F$ stays silent. $H_0$ commits; timeout recovery later re-certifies the
old chain as a timerless ANCHOR. Repeating this execution can leave $H_1$
without any non-genesis commitment. There is no equivocation to exclude.
\Cref{alg:steady} prevents this execution by processing unseen history
proposals and preserving earlier-term voting eligibility.

\paragraph{B. Shorten normal-term deadlines to $2p\Delta$.}
Use the same committee and $K=1$. Let $H_0$ propose at $s$ and self-vote.
Choose $0<\varepsilon<\Delta/3$. Fast exchange with $F$ gives $H_0$ the
terminal QC and successor entry at $x=s+2\varepsilon$, while $H_1$ first
receives the proposal and votes at $s+\Delta$. $F$, the successor, sends
no DATA proposal and makes its BLAME share available to $H_0$ in advance.
$H_0$ commits at $s+2\Delta$, then under the shortened rule sends BLAME
and freezes at $x+2\Delta$. Deliver the resulting proof to $H_1$ within
$\varepsilon$. It arrives at
\[
 s+2\Delta+3\varepsilon<s+3\Delta,
\]
canceling $H_1$'s timer before its expiry. All delays obey the upper bound.
Recovery is timerless on the anchor. Repeat with $H_0$ as initializer:
$H_0$ commits increasing heights but $H_1$ never commits a non-genesis
block. $F$ only withholds, so exclusion does not stop the attack.
Thus a deadline sufficient to avoid blaming a correct successor need not
preserve all-replica liveness. The specified $6\Delta$ first check leaves
the $3\Delta$ window of \cref{lem:initializer} intact.

\paragraph{C. Remove permanent exclusion.}
For any fixed finite $K$, let a correct initializer and $F$ exchange votes
with a sufficiently small positive delay $\varepsilon$. They can certify
all $K$ blocks and admit $F$ as successor before the first $2\Delta$
timer expires. $F$ signs two different valid first children of the terminal
block. Deliver their equivocation proof to all correct replicas before
any timer expires. They cancel every timer and recover. Repeat in every
generation with the same still-eligible $F$. QCs and heights can grow
while no non-genesis block commits. A finite quota, including $K=n$,
does not exclude this execution when the model supplies no lower bound
that forces a long enough proposal window. After the first revelation
identifies $g_{\mathrm{bad}}$, the specified finite $G_p$ prohibits reusing
the same leader in every later generation; the execution removes that condition.

\paragraph{D. Start the handoff bound at private QC formation.}
Let $F$ lead a $K=1$ term, with correct successor $H_0$ and another correct
replica $H_1$, and choose $0<\varepsilon<\Delta$.
$H_0$ receives the terminal proposal and votes at $s$.
$F$ privately combines that vote with its own at $q=s+\varepsilon$,
but withholds its vote and QC. Deliver $H_0$'s forwarded proposal and vote
to $H_1$ at $s+\Delta$, and $H_1$'s vote and backup QC to $H_0$ at
$s+2\Delta$.
$H_0$ cannot obtain the terminal QC before $s+2\Delta>q+\Delta$.
The bound in \cref{lem:handoff} instead starts at a correct replica's
complete verified QC receipt, which triggers actual forwarding.

\paragraph{E. Treat two unauthenticated messages as equivocation.}
An outsider can fabricate different blocks for the same position and
force arbitrary freezes if signatures and admission are not checked.
Evidence-first processing therefore prioritizes verification of a proof;
it does not bypass that verification. Conversely, a malformed enclosing
message must not suppress a valid proof already present in its signed
components.

\section{Communication and latency}
\label{sec:complexity}
Count correct-generated complete admitted protocol traffic, excluding
arbitrary Byzantine flooding and induced retrieval traffic as specified below.
Each correct replica forwards and votes for at most one DATA
block per position, sends at most one BLAME share per term, and broadcasts
its first quit proof once per generation. Let $b$ bound a fresh payload,
$\sigma$ a signature, and $A$ a bound on additional dependency bytes in
each normalized bundle charged to the operation. Different operations may
have different $A$. Explicit QCs have $O(n\sigma)$ bytes.

\begin{center}\small
\begin{tabularx}{\linewidth}{@{}>{\raggedright\arraybackslash}p{.24\linewidth}
 >{\raggedright\arraybackslash}p{.20\linewidth}Y@{}}
\toprule
Operation & Correct messages & Conservative bytes\\
\midrule
DATA forwarding and votes & $O(n^2)$ per position &
$O(n^2(b+A+n\sigma))$\\
Terminal-QC backup & $O(n)$ per certified terminal &
$O(n(b+A+n\sigma))$ to one successor\\
Progress complaints & $O(n^2)$ per admitted term &
$O(n^2\sigma)$ plus admission/context data\\
Quit and recovery & $O(n^2)$ per recovery &
$O(n^2(A+n\sigma))$ for proof, STATUS, NEWVIEW, and ANCHOR votes\\
Local proposal timers & No messages & No extra vote or QC phase\\
L1 bulletin board & External service & Evidence submission, public authority
reads, and eligibility updates\\
\bottomrule
\end{tabularx}
\end{center}
Different honest replicas can initially accept different equivocations,
but each still sends only one DATA vote/forward per position before proof
processing freezes the generation. Certificate subsets do not defeat
semantic deduplication. Let $S$ count distinct processed DATA positions
$(g,t,j)$, $T$ distinct admitted terms $(g,t)$ including unfinished ones,
and $R$ generations whose recovery has been initiated, including those
still waiting for L1. The message bound is
\[
 O\bigl(n^2(S+T+R)+nT\bigr).
\]
Partial BLAME traffic is charged even when no recovery follows. For
normally completed $K$-block terms, backup forwarding costs $O(n/K)$
messages per block; complaints cost at most $O(n^2/K)$ per block.
Votes and old DATA processed during recovery remain charged to their
originating positions. The bulletin-board model requires no BFT relay of
L1 finality proofs; its service costs are separate from these message counts.

\paragraph{History and retrieval costs.}
The displayed bounds cover complete bundles. Malformed or incomplete
Byzantine traffic can induce arbitrary requests; it is outside this
bounded-traffic accounting. Any implementation that sends cache references
must separately count requests, replies, and transferred missing objects,
and preserve the complete-message timing assumption. $A$ is not constant:
without cache reuse, forwarding length-$H$ prefixes can require
$A=O(H(b+n\sigma))$ plus control evidence. Across $S$ successive positions,
repeated full-prefix transmission can cost $O(n^2S^2(b+n\sigma))$ bytes,
even though DATA message counts remain $O(n^2S)$. No linear amortized
history-byte claim or constant-size QC assumption is made.

\paragraph{Normal operation.}
A proposal accepted at $\tau_r(P)$ commits locally at
$\tau_r(P)+2\Delta$ if its timer survives. From a correct vote at $s$,
an uninterrupted window gives commitment everywhere by $s+3\Delta$.
From the first correct terminal-QC time $q$, a correct successor proposes
by $q+\Delta$, all receive by $q+2\Delta$, and surviving timers commit
that child everywhere by $q+4\Delta$. With homogeneous actual delay
$\delta$, ordinary proposal/QC iterations take about $2\delta$.
The backup adds no mandatory $\Delta$ wait: a leader acts as soon as
its parent QC is available. It does add the stated $O(n)$ terminal traffic.

\Needspace{13\baselineskip}
\paragraph{Recovery.}
Let $a=\min_r a_r$ be the first correct quit time. The common L1
availability time gives
\[
 e=\min_r e_r(g+1,0)=\max\{a+\Delta,d(g+1)\}.
\]
With a correct initializer, the first DATA proposal reaches all by
$e+6\Delta$ and an uninterrupted window commits it everywhere by
\[
 \max\{a+\Delta,d(g+1)\}+8\Delta.
\]
If $d(g+1)\le a+\Delta$, this is $a+9\Delta$; otherwise recovery
waits for the public authority. Eventual exclusion and
authority availability give no numerical bound on the last equivocation,
number of aborted generations, or end-to-end recovery delay. Quantitative
bounds require corresponding L1 and scheduling bounds.



\begin{thebibliography}{9}
\bibitem{sync}
I. Abraham, D. Malkhi, K. Nayak, L. Ren, and M. Yin.
\emph{Sync HotStuff: Simple and Practical Synchronous State Machine Replication}.
IEEE Symposium on Security and Privacy, 2020, pp. 106--118.
\href{https://conferences.computer.org/sp/pdfs/sp/2020/349700a106.pdf}{Published paper}.
\bibitem{schneider}
F. B. Schneider. Implementing fault-tolerant services using the state machine
approach: a tutorial. \emph{ACM Computing Surveys}, 22(4):299--319, 1990.
\href{https://doi.org/10.1145/98163.98167}{DOI}.
\end{thebibliography}
\end{document}
